\subsubsection{Simplexný a duplexný reťazec}
\begin{itemize}
    \item \textbf{Simplex (ssDNA)} --- jednoreťazcová DNA. Obsahuje iba jedno vlákno nukleotidov. Vyskytuje sa napríklad počas replikácie alebo pri niektorých vírusoch. V kontexte nanopórového sekvenovania sa čítanie vykonáva práve na jednoreťazcovej forme.
    \item \textbf{Duplex (dsDNA)} --- dvojreťazcová DNA. Dve komplementárne vlákna sú navzájom spojené vodíkovými väzbami medzi pármi báz (A--T, G--C) a tvoria charakteristickú dvojitú špirálu (double helix). Toto je štandardná forma, v ktorej je DNA uložená v bunke.
\end{itemize}

Pri biosyntéze DNA vzniká z jednoreťazcového templátu nový komplementárny reťazec --- výsledkom je duplexná DNA. Pre každý pôvodný reťazec tak existuje komplementárny reťazec s opačnou polaritou. To znamená, že pri nanopórovom sekvenovaní môžeme naraziť na štyri varianty toho istého úseku DNA:
\begin{enumerate}
    \item pôvodný reťazec v smere 5'$\rightarrow$3',
    \item pôvodný reťazec v smere 3'$\rightarrow$5' (reverzný),
    \item komplementárny reťazec v smere 5'$\rightarrow$3',
    \item komplementárny reťazec v smere 3'$\rightarrow$5' (reverzný komplement).
\end{enumerate}

Tieto štyri varianty sú kľúčovým zdrojom komplikácií pri klastrovacích algoritmoch --- sekvencie reprezentujúce rovnaký úsek DNA môžu byť sekvenerom prečítané v rôznych formách a bez primerov ich nie je možné jednoducho stotožniť.
